% ============================================================
%  Georg Brandes, Levned, Bind I
%  Kjøbenhavn: Gyldendal, 1905.
%  Udvalgte passager om Rasmus Nielsen og Hans Brøchner
%  TRANSCRIPTION (Danish)
%  Source: https://georgbrandes.dk (Det Danske Sprog- og
%          Litteraturselskab, Digitale Hovedstrømninger)
% ============================================================
\documentclass[12pt,a4paper]{article}
\usepackage[utf8]{inputenc}
\usepackage[T1]{fontenc}
\usepackage{libertinus}
\usepackage{libertinust1math}
\usepackage[danish]{babel}
\usepackage[protrusion=true,expansion=true]{microtype}
\usepackage{setspace}
\usepackage[margin=1.2in]{geometry}
\usepackage{fancyhdr}
\usepackage{hyperref}

\hypersetup{
  pdftitle={Levned I — Passager om R. Nielsen og H. Brøchner},
  pdfauthor={Georg Brandes},
  hidelinks
}

\pagestyle{fancy}
\fancyhf{}
\rhead{Brandes, \emph{Levned} I (1905)}
\cfoot{\thepage}

\setstretch{1.3}

% Original page-number marker
\newcommand{\pnum}[1]{\noindent\makebox[0pt][r]{\small\textit{[#1]}\quad}}

\title{\textbf{Levned}, Bind I\\[6pt]
  \large Udvalgte passager om Rasmus Nielsen og Hans Brøchner}
\author{Georg Brandes}
\date{Kjøbenhavn: Gyldendal, 1905}

\begin{document}
\maketitle
\thispagestyle{fancy}

\bigskip

\noindent\textit{Nedenstående passager er udvalgt for deres relevans for
forholdet mellem Georg Brandes og hans to filosofiprofessorer, Rasmus
Nielsen og Hans Brøchner, samt for Tro-og-Viden-striden i dansk
1860'er-filosofi.}

\bigskip\hrule\bigskip

\subsection*{Kapitel 3 (\emph{Overgangsaar}), §~5}

\noindent Kontekst: En udenlandsk student ved navn Kappers, af vestindisk
oprindelse, studerede ved Københavns Universitet og forberedte sig til den
filosofiske eksamen.

\medskip

\pnum{85}Kappers tilbragte næsten hele sin Dag drivende om paa Gaderne
i Samtale med Kammerater; altid var han rede til en Spadseretur; aldrig
saa man ham arbejde eller hørte man ham tale om sit Arbejde. Ikke
destomindre indstillede han, Udlændingen, der lige akkurat var Sproget
mægtig, sig efter kun et halvt Aars Forløb --- uden altsaa at have hørt
Forelæsningerne til Ende --- til den filosofiske Eksamen og tog den med
Udmærkelse i alle tre Fag; ja Rasmus Nielsen, som eksaminerede ham i
Propædeutik, blev saa begejstret over denne Fremmedes hurtige og kloge
Svar, at han gav ham \textit{Særdeles udmærket godt}, en Karakter, som
ikke eksisterte.

\bigskip\hrule\bigskip

\subsection*{Kapitel 4 (\emph{Ynglingeaar}), §~2}

\noindent Kontekst: Brandes er netop blevet student (efteråret 1859).
Brøchner havde tilbudt filosofisk vejledning til studerende. Brandes
beskriver sit første besøg hos ham den 2.\ september 1861.

\medskip

Længe havde jeg drevet mine ikke-juridiske Studier, som jeg kunde bedst
uden Vejledning af nogensomhelst Lærer. Men jeg havde følt Savnet af en
saadan. Og da en nyansat Docent ved Universitetet, Professor H. Brøchner,
tilbød Vejledning i Studiet af Filosofi til hvem der vilde indfinde sig
hos ham paa visse opgivne Tider, havde jeg følt mig stærkt fristet til at
benytte mig deraf. Jeg havde havt Betænkeligheder derved, thi jeg vilde
nødigt opgive det mindste af min dyrebare Frihed; jeg nød min tilbagetrukne
Tilstand, mit beskedne Studielivs Hemmelighedsfuldhed; men paa den anden
Side vilde jeg ikke kunne give mig i Kast med Plato og Aristoteles uden
Vink af en kyndig Haand om hvad og hvorledes.

Jeg var da i stor Spænding. Jeg havde hørt Professor Brøchner tale om
Psykologi; men Foredraget var behandlet med en saa ængstelig Omhu, var
saa ensformigt og fremmedlydende, at det ikke kunde andet end skræmme.
Alene Professorens fornemme Skønhed tiltrak, især de store sorgfulde Øjne
med det dejlige Blik, der saa ud som havde de forsket og grædt.

\pnum{130}Nu besluttede jeg at vove mig op til Brøchner. Men jeg
havde ikke Mod til forud at tale derom til min Moder for ikke at
forskrække mig selv ved at nævne det; saa urolig var jeg ved Skridtet.
Da Dagen kom, jeg havde bestemt til Besøget --- det var den 2.\ September
1861 --- gik jeg flere Gange op og ned foran Huset, førend jeg kunde
bestemme mig til at gaa op ad Trappen; jeg søgte forud at beregne, hvad
Professoren vel vilde sige og hvad jeg da rettest skulde svare, i det
Uendelige.

Den høje, smukke Mand, der lignede en spansk Ridder, lukkede selv Døren
op og modtog venligt det unge Menneske, der snart skulde blive hans
nærmeste Elev. Han vidste til min Forundring, straks han hørte mit Navn,
fra hvilken Skole jeg var udgaaet og at jeg nylig var bleven Student. Han
fraraadede energisk en Gennemgaaen af Plato og Aristoteles som altfor
anstrengende --- sagde eller antydede, at jeg burde forskaanes for den
besværlige Vej, han selv havde tilbagelagt, og skitserede en Studieplan
for nyere Filosofi og Æstetik. Hans Holdning indgav Mod og efterlod det
Hovedindtryk, at han vilde spare Begynderen enhver unyttig Anstrengelse.
Jeg følte mig med mine unge Kræfter lige fuldt, da jeg gik hjem, en
Smule skuffet ved ikke at skulle tage fat paa Filosofiens Historie forfra.

Besøget gentoges snart, og hurtigt udspandt der sig mellem Lærer og Elev
det inderligste Forhold, fra den ældres Side et faderligt Velviljesforhold,
som den yngre endnu ikke havde kendt, under hvilket Læreren, undervisende
og praktisk vejledende, uafbrudt havde Opmærksomheden henvendt paa Alt,
hvad der kunde fremme Lærlingens Vel og sikre hans Fremtid; fra den yngres
Side et Ærbødigheds- og Hengivenhedsforhold, en dyb Taknemmelighed for den
Omsorg, hvis Genstand han var.

Jeg følte vel nu og da overfor denne Lærers store \pnum{131}Grundighed
og hans Dygtighed til at tumle de vanskeligste Tanker en pinlig Mistillid
til mine egne Evner og min egen Aandskraft i Sammenligning med hans. Jeg
følte mig heller ikke sjældent ilde berørt ved, at der ligesom kom noget
Skævt ind i vort Forhold, naar Brøchner tiltrods for de Tvivl og
Indvendinger, jeg fremsatte, altid tog som givet, at jeg maatte dele hans
panteistiske Grundanskuelse, uden Blik for, at jeg endnu omtumledes i
Tvivl og famlede efter et Fodfæste. Men det fortrolige Forhold til den
modnede Mand havde en Tiltrækning, som Forholdet til de usikre eller
ungdommeligt indskrænkede Kammerater nødvendigvis maatte mangle: han
havde et Livs Erfaring bag sig, han saa fra oven nedad paa et ungt
Menneskes Sympatier og Antipatier.

[\ldots]

Blot den Maade, paa hvilken Brøchner en Dag med et Smil bemærkede:
De læser nok af Princip ikke \textit{Dagbladet}, fik mig til i et
Glimt at indse det Komiske i min Harme over en Hjerrilds Artikler.

[\ldots]

\pnum{132}I væsenlig Overensstemmelse med Hans Brøchner følte jeg
mig først flere Aar efter Bekendtskabets Stiftelse. Men da havde et
Studium af Ludwig Feuerbachs Skrifter henrevet mig; thi Feuerbach var
den første Tænker, hos hvem jeg fandt Gudsideens Opstaaen i det
menneskelige Sind fyldestgørende forklaret. Hos Feuerbach mødte der mig
endelig en Fremstilling uden Omsvøb og uden den tyske Filosofis sædvanlige
tunge Formler, en sejrrig Klarhed, som virkede højst velgørende paa mit
eget Tankesæt og gav mig Tryghed. Havde jeg i flere Aar følt mig
betydeligt tilhøjre for min Ven og Lærer, saa kom der nu en Tid, hvor
jeg følte mig adskilligt tilvenstre for ham med hans hemmelighedsfulde
Leven i det evige Aandsrige, af hvilket han følte sig som Led.

\bigskip\hrule\bigskip

\subsection*{Kapitel 4 (\emph{Ynglingeaar}), §~5}

\noindent Kontekst: Efter at have modtaget et accessit (men ikke
guldmedaljen) for sin prisafhandling opsøger Brandes sine tre
eksamenskommissærer --- Hauch, Sibbern og Brøchner --- for at
præsentere sig. Her er uddrag fra besøget hos Sibbern.

\medskip

\pnum{137}Sibbern indlod sig paa Enkeltheder i Afhandlingen og var
strengere end Hauch. Han ankede over, at der fattedes det egenlige Afsnit;
Præmisserne var altfor vidtløftigt affattede. Han raadede til et mere
historisk, mindre filosofisk Studium af Literatur og Kunst. Han var
tilfreds ved at høre om mit nære Forhold til Brøchner, medens Hauch
hellere havde set mig knyttet til Rasmus Nielsen, hvem han skemtende
betegnede som »en rigtig Rugbrødsnatur«.

\bigskip\hrule\bigskip

\subsection*{Kapitel 4 (\emph{Ynglingeaar}), §§~11--13}

\noindent Kontekst: Efteråret 1862 og begyndelsen af 1863. Brandes sender
sin store afhandling om \textit{Romeo og Julie} til Brøchner og modtager
kort efter et afgørende nytårsbrev.

\medskip

\pnum{153}[\ldots] jeg udarbejdede i de sidste Maaneder af 1862 en
meget stor Afhandling om Shakespeares \textit{Romeo og Julie}, der
væsenlig drejede sig om Tragediens Grundproblemer, som de opfattedes i
Datidens Æstetik; den er gaaet tabt, som saa meget Andet fra disse Aar.
Jeg sendte den til Professor Brøchner og udbad mig dennes Dom.

[\ldots]

\pnum{154}Første Januar 1863 modtog jeg et Nytaarsbrev fra Brøchner,
hvori denne skrev, at Afhandlingen om \textit{Romeo og Julie} havde gjort
et stærkt Indtryk paa ham, og at efter hans Mening kunde umuligt Nogen
gøre mig Rangen stridig med Hensyn til det Professorat i Æstetik, der
efter Sagens Natur snart maatte blive ledigt, da Hauch i sin høje Alder
ikke kunde beklæde Lærestolen ret længe.

Saaledes indlededes af min ivrige Beskytter det en Menneskealder igennem
fortsatte Vrøvl om mit Forhold til et Embede, der efterhaanden i mit Liv
blev hvad Franskmændene kalder \textit{une scie}, et irriterende
Veksérspil, hvori jeg selv ikke havde ringeste Del, men som klæbede ved
mit Navn.

\pnum{155}Hint Brev satte mig i Sindsbevægelse; ikke fordi der i saa
ung en Alder stilledes mig en ærefuld Samfundsstilling i Udsigt af en
Mand, der havde Sagkundskab til Bedømmelse af min Adkomst til den, men
fordi denne smilende Udsigt til et Embede i mine Øjne var en Fælde, der
kunde fange mig saaledes, at jeg ikke kom til at gaa ad den Forsagelsens
Vej, der ene syntes mig at føre til mit Livsmaal. [\ldots] Det var
vanskeligere at nedlægge et Professorat end aldrig at overtage det. Og
var man først Professor, blev man snart gift og bosat Borger i Staten,
ude af Stand til at vove. At anlægge sit Liv efter Brøchners Opfordring
var som at sælge sin Sjæl til Fanden.

Jeg svarte da kortelig, at jeg holdt for meget af Hauch til at kunne
eller ville spekulere i hans Død. Dog hertil indvendte Brøchner meget
logisk: Jeg spekulerer ikke i hans Død, men i hans Liv, thi jo længer han
lever, des bedre er De forberedt til at blive hans Efterfølger.

\bigskip\hrule\bigskip

\subsection*{Kapitel 4 (\emph{Ynglingeaar}), §~13}

\noindent Kontekst: Slutningen af 1863. Danmark er i krig; Brandes
forbereder sig til sin magistereksamen og skriver til Brøchner.

\medskip

\pnum{161}[\ldots] Jeg skrev i Slutningen af 1863 til min Lærer,
Professor Brøchner, der havde lovet mig en kort filosofisk Oversigt som
Forberedelse til Universitetsprøven: »Jeg indstiller mig under de aller
ugunstigste Omstændigheder, eftersom mit Hoved nu i mere end en Maaned
har været saa optaget af Dagens Begivenheder, at jeg hverken har kunnet
læse eller tænke, og selve Eksamenstiden synes at ville blive endnu
uroligere. Dog tør jeg nu ikke vige tilbage, da jeg ellers risikerer ---
hvad mulig alligevel sker --- at blive forhindret fra Eksamen ved Session
og Indkaldelse, og saaledes maaske kommer til at ligge i min Grav som
\textit{studiosus} istedenfor som \textit{candidatus magisterii}, hvilket
sidste dog tager sig ulige stateligere ud og er mere tilfredsstillende
for et ærekært Menneske som Deres ærbødige og inderligt hengivne osv.«

\bigskip\hrule\bigskip

\subsection*{Kapitel 4 (\emph{Ynglingeaar}), §~14}

\noindent Kontekst: Kort inden eksamen i begyndelsen af 1864 besøger
Brandes for første gang Rasmus Nielsen. (Teksten er afbrudt i den
tilgængelige kilde ved siden 162.)

\medskip

Kort forinden havde jeg for første Gang besøgt Professor Rasmus Nielsen.
Denne var højst forekommende, kendte mig, som han maaske huskede at have
eksamineret [\ldots]

\bigskip\hrule\bigskip

\subsection*{Kapitel 6 (\emph{Ungdomsliv}), §~1}

\noindent Kontekst: Brandes er vendt hjem fra sit første lange
Frankrig-ophold (1866--67) og kaster sig ind i Tro-og-Viden-striden.

\medskip

\pnum{203}Efter Tilbagekomsten fra Frankrig til Danmark i 1867
optoges mine Tanker paany af den Fejde, som var brudt ud i dansk
Literatur mellem Videnskab og saakaldt Aabenbaring (i Datidens Sprog:
Tro og Viden). Flere og flere havde efterhaanden udgivet Indlæg i denne
Strid, og jeg, der havde min højeste aandelige Interesse i den, tog Del i
Ordskiftet med dobbelt Front, dels imod de rettroende Teologer, dels og
især imod R.\ Nielsen, Teologernes Angriber, der i mine Øjne ikke var
mindre teologisk sindet end sine Modstandere.

\bigskip\hrule\bigskip

\subsection*{Kapitel 6 (\emph{Ungdomsliv}), §~19}

\noindent Kontekst: Brøchners personlige elever i slutningen af 1860'erne.

\medskip

Iblandt Brøchners personlige Elever var en ung Studerende ved Navn
Kristian Møller, der udelukkende gav sig af med Filosofi og som var
Brøchner ganske særligt kær. \pnum{239}Det var et ualmindelig skarpt
Hoved, anlagt til kritisk opløsende Forskning. Han gav forsaavidt store
Løfter, som han syntes at maatte kunne begrunde tilintetgørende Domme
over meget, der var taaget eller selvmodsigende og dog anset; men han
var iøvrigt en mærkeligt ufrugtbar Aand. Han havde ingen Indfald, ingen
Fantasi. Man kunde ikke vente, at han vilde komme til at frembringe
synderligt nye Tanker, men at han vilde gennemskue Uklarhed, gennemtænke
Vildfarelser og med Aarene yde et grundigt Rydningsarbejde.

\bigskip\hrule\bigskip

\subsection*{Kapitel 7 (\emph{Andet længere Udenrigsophold}), §~1}

\noindent Kontekst: Brandes er i Hamborg på vej til Paris i 1867.
Han besøger zoologisk have og reflekterer over naturen og dens
modsætning til Nielsens filosofiske terminologi.

\medskip

\pnum{270}[\ldots] Og det morede mig, at R.\ Nielsen og hans Elever
stadigt talte, som stod de paa den fortroligste Fod med selve »det
Centrale« i Tilværelsen, og haanede Andre, der formentlig ikke formaaede
at komme i Forhold dertil.

\end{document}
%%% Local Variables:
%%% mode: latex
%%% TeX-master: t
%%% End:
